\section{Reduced preschemes; separation condition}
\label{section:I.5}

\subsection{Reduced preschemes.}
\label{subsection:I.5.1}

\begin{proposition}[5.1.1]
\label{I.5.1.1}
Let $(X,\mathscr{O}_X)$ be a prescheme and $\mathscr{B}$ a quasi-coherent
$\mathscr{O}_X$-algebra. There exists one and only one quasi-coherent
$\mathscr{O}_X$-module $\mathscr{N}$ whose stalk
$\mathscr{N}_x$ at every $x\in X$ is the nilradical of the ring
$\mathscr{B}_x$. When $X$ is affine, and consequently
$\mathscr{B}=\widetilde{B}$, where $B$ is an algebra over $A(X)$, we have
$\mathscr{N}=\widetilde{\mathfrak{N}}$, where $\mathfrak{N}$ is the nilradical
of $B$.
\end{proposition}

\oldpage[I]{128}
The question being local, we are reduced to proving the last assertion. We
know that $\widetilde{\mathfrak{N}}$ is a quasi-coherent $\mathscr{O}_X$-module
(\hyperref[I.1.4.1]{1.4.1}) and that its stalk at the point
$x\in X$ is the ideal $\mathfrak{N}_x$ of the ring of fractions $B_x$; it all
comes down to proving that the nilradical of $B_x$ is contained in
$\mathfrak{N}_x$, the opposite inclusion being evident. Now let $z/s$ be an
element of the nilradical of $B_x$, with $z\in B$, $s\notin\mathfrak{j}_x$; by
hypothesis, there exists an integer $k$ such that $(z/s)^k=0$, which means
that there exists $t\notin\mathfrak{j}_x$ such that $tz^k=0$. We conclude that
$(tz)^k=0$, and consequently $z/s=(tz)/(ts)$ belongs to $\mathfrak{N}_x$.

We will say that the quasi-coherent $\mathscr{O}_X$-module $\mathscr{N}$ thus
defined is the \emph{nilradical} of the $\mathscr{O}_X$-algebra $\mathscr{B}$;
in particular, we will denote by $\mathscr{N}_X$ the nilradical of
$\mathscr{O}_X$.

\begin{corollary}[5.1.2]
\label{I.5.1.2}
Let $X$ be a prescheme; the closed sub-prescheme of $X$ defined by the sheaf
of ideals $\mathscr{N}_X$ is the only reduced sub-prescheme
(\hyperref[0.4.1.4]{0, 4.1.4}) of $X$ having $X$ as underlying space;
it is also the smallest sub-prescheme of $X$ having $X$ as underlying
space.
\end{corollary}

As the structural sheaf of the closed sub-prescheme $Y$ defined by
$\mathscr{N}_X$ is $\mathscr{O}_X/\mathscr{N}_X$, it is immediate that $Y$
is reduced and has $X$ as underlying space, since
$\mathscr{N}_x\ne\mathscr{O}_x$ for every $x\in X$. To prove the other
assertions, note that a sub-prescheme $Z$ of $X$ having $X$ as underlying
space is defined by a sheaf of ideals $\mathscr{J}$
(\hyperref[I.4.1.3]{4.1.3}) such that
$\mathscr{J}_x\ne\mathscr{O}_x$ for every $x\in X$. We may restrict to the case
where $X$ is affine, say $X=\operatorname{Spec}(A)$ and
$\mathscr{J}=\widetilde{\mathfrak{J}}$, where $\mathfrak{J}$ is an ideal of
$A$; then, for every $x\in X$, we have
$\mathfrak{J}\subset\mathfrak{j}_x$, so $\mathfrak{J}$ is contained in
all the prime ideals of $A$, that is, in their intersection
$\mathfrak{N}$, the nilradical of $A$. This proves that $Y$ is the smallest
sub-prescheme of $X$ having $X$ as underlying space
(\hyperref[I.4.1.9]{4.1.9}); moreover, if $Z$ is distinct from $Y$, we
necessarily have $\mathscr{J}_x\ne\mathscr{N}_x$ for at least one $x\in X$,
and consequently (\hyperref[I.5.1.1]{5.1.1}) $Z$ is not reduced.

\begin{definition}[5.1.3]
\label{I.5.1.3}
We call the reduced prescheme associated to a prescheme $X$, and denote by
$X_{\mathrm{red}}$, the unique reduced sub-prescheme of $X$ having $X$ as
underlying space.
\end{definition}

To say that a prescheme $X$ is reduced therefore means that
$X=X_{\mathrm{red}}$.

\begin{proposition}[5.1.4]
\label{I.5.1.4}
For the prime spectrum of a ring $A$ to be a reduced (resp.
integral) prescheme (\hyperref[I.2.1.8]{2.1.8}), it is necessary and
sufficient that $A$ be a reduced (resp. integral) ring.
\end{proposition}

Indeed, it follows immediately from (\hyperref[I.5.1.1]{5.1.1}) that the
condition $\mathscr{N}=(0)$ is necessary and sufficient for
$X=\operatorname{Spec}(A)$ to be reduced; the assertion concerning integral
rings is then a consequence of
(\hyperref[I.1.1.13]{1.1.13}).

Since every ring of fractions $\ne\{0\}$ of an integral ring is integral, it
follows from (\hyperref[I.5.1.4]{5.1.4}) that for every \emph{locally
integral} prescheme $X$, $\mathscr{O}_x$ is an \emph{integral} ring
for every $x\in X$. The converse is true when the underlying space of
$X$ is \emph{locally Noetherian}: indeed $X$ is then reduced, and if
$U$ is an affine open subset of $X$ which is a Noetherian space, $U$ has only
finitely many irreducible components, so its ring $A$ has only
finitely many minimal prime ideals
(\hyperref[I.1.1.14]{1.1.14}). If two of these components $U_i$ had a
common point $x$, $\mathscr{O}_x$ would have at least two distinct minimal
prime ideals, and so would not be integral; consequently the $U_i$ are
pairwise disjoint open subsets, and each of them is therefore integral.

\begin{env}[5.1.5]
\label{I.5.1.5}
Let $f=(\psi,\theta):X\to Y$ be a morphism of preschemes; the homomor-
\oldpage[I]{129}
phism $\theta_x^\sharp:\mathscr{O}_{\psi(x)}\to\mathscr{O}_x$ maps every
nilpotent element of $\mathscr{O}_{\psi(x)}$ to a nilpotent element of
$\mathscr{O}_x$; by passing to quotients, we therefore deduce from
$\theta^\sharp$ a homomorphism
\[
  \omega:\psi^*(\mathscr{O}_Y/\mathscr{N}_Y)
  \longrightarrow \mathscr{O}_X/\mathscr{N}_X
\]
it is clear that for every $x\in X$,
$\omega_x:\mathscr{O}_{\psi(x)}/\mathscr{N}_{\psi(x)}
\to\mathscr{O}_x/\mathscr{N}_x$ is a local homomorphism, so
$(\psi,\omega^b)$ is a morphism of preschemes
$X_{\mathrm{red}}\to Y_{\mathrm{red}}$, which we will denote by
$f_{\mathrm{red}}$ and call the \emph{reduced} morphism associated to $f$. It
is immediate that for two morphisms $f:X\to Y$, $g:Y\to Z$, we have
$(g\circ f)_{\mathrm{red}}=g_{\mathrm{red}}\circ f_{\mathrm{red}}$, so we have
defined $X_{\mathrm{red}}$ as a \emph{covariant functor} in $X$.

The preceding definition shows that the diagram
\[
  \xymatrix{
    X_{\mathrm{red}}\ar[r]^{f_{\mathrm{red}}}\ar[d] &
    Y_{\mathrm{red}}\ar[d]\\
    X\ar[r]_f & Y
  }
\]
is commutative, the vertical arrows being the injection morphisms; in
other words, $X_{\mathrm{red}}\to X$ is a \emph{functorial} morphism.
Note in particular that if $X$ is \emph{reduced}, every morphism
$f:X\to Y$ factors as
$X\xrightarrow{f_{\mathrm{red}}}Y_{\mathrm{red}}\to Y$; in other words,
$f$ is \emph{dominated} by the injection morphism
$Y_{\mathrm{red}}\to Y$.
\end{env}

\begin{proposition}[5.1.6]
\label{I.5.1.6}
Let $f:X\to Y$ be a morphism; if $f$ is surjective (resp. radicial, an
immersion, a closed immersion, an open immersion, a local immersion,
a local isomorphism), then so is $f_{\mathrm{red}}$. Conversely,
if $f_{\mathrm{red}}$ is surjective (resp. radicial), then so is
$f$.
\end{proposition}

The proposition is trivial if $f$ is surjective; if $f$ is radicial, it
follows from the fact that for every $x\in X$, the field $k(x)$ is the same
for the preschemes $X$ and $X_{\mathrm{red}}$
(\hyperref[I.3.5.8]{3.5.8}). Finally, if $f=(\psi,\theta)$ is an
immersion, a closed immersion or a local immersion (resp. an open immersion,
or a local isomorphism), the proposition follows from the fact that if
$\theta_x^\sharp$ is surjective (resp. bijective), then so is the
homomorphism obtained by passing to the quotients by the nilradicals of
$\mathscr{O}_{\psi(x)}$ and $\mathscr{O}_x$
(\hyperref[I.5.1.2]{5.1.2} and \hyperref[I.4.2.2]{4.2.2})
(cf. (\hyperref[I.5.5.12]{5.5.12})).

\begin{proposition}[5.1.7]
\label{I.5.1.7}
If $X$, $Y$ are two $S$-preschemes, the preschemes
$X_{\mathrm{red}}\times_{S_{\mathrm{red}}}Y_{\mathrm{red}}$ and
$X_{\mathrm{red}}\times_S Y_{\mathrm{red}}$ are identical, and
are canonically identified with a sub-prescheme of $X\times_S Y$ having the
same underlying space as this product.
\end{proposition}

The canonical identification of $X_{\mathrm{red}}\times_S Y_{\mathrm{red}}$
with a sub-prescheme of $X\times_S Y$ having the same underlying space follows
from (\hyperref[I.4.3.1]{4.3.1}). On the other hand, if $\varphi$ and
$\psi$ are the structural morphisms $X_{\mathrm{red}}\to S$,
$Y_{\mathrm{red}}\to S$, they factor through $S_{\mathrm{red}}$
(\hyperref[I.5.1.5]{5.1.5}), and since $S_{\mathrm{red}}\to S$ is a
monomorphism, the first assertion follows from
(\hyperref[I.3.2.4]{3.2.4}).

\begin{corollary}[5.1.8]
\label{I.5.1.8}
The preschemes $(X\times_S Y)_{\mathrm{red}}$ and
$(X_{\mathrm{red}}\times_{S_{\mathrm{red}}}Y_{\mathrm{red}})_{\mathrm{red}}$
are canonically identified.
\end{corollary}

This follows from (\hyperref[I.5.1.2]{5.1.2} and
\hyperref[I.5.1.7]{5.1.7}).

Note that if $X$ and $Y$ are reduced preschemes, the same is not
necessarily true of $X\times_S Y$, because the tensor product of two
reduced algebras can have nilpotent elements.

\begin{proposition}[5.1.9]
\label{I.5.1.9}
\oldpage[I]{130}
Let $X$ be a prescheme and $\mathscr{J}$ a quasi-coherent sheaf of ideals of
$\mathscr{O}_X$ such that $\mathscr{J}^n=0$ for some integer $n>0$. Let $X_0$
be the closed sub-prescheme $(X,\mathscr{O}_X/\mathscr{J})$ of $X$; for $X$
to be an affine scheme, it is necessary and sufficient that $X_0$ be one.
\end{proposition}

The condition being evidently necessary, let us prove that it is sufficient.
If we set $X_k=(X,\mathscr{O}_X/\mathscr{J}^{k+1})$, it all comes down to
proving by induction on $k$ that the $X_k$ are affine, so we are reduced to
the case where $\mathscr{J}^2=0$. Set
\[
  A=\Gamma(X,\mathscr{O}_X),\qquad
  A_0=\Gamma(X_0,\mathscr{O}_{X_0})
      =\Gamma(X,\mathscr{O}_X/\mathscr{J}).
\]
From the canonical homomorphism
$\mathscr{O}_X\to\mathscr{O}_X/\mathscr{J}$ we deduce a ring homomorphism
$\varphi:A\to A_0$. We will see below that $\varphi$ is
\emph{surjective}, so that the sequence
\[
  0\to\Gamma(X,\mathscr{J})\to\Gamma(X,\mathscr{O}_X)
  \to\Gamma(X,\mathscr{O}_X/\mathscr{J})\to 0
  \tag{5.1.9.1}\label{I.5.1.9.1}
\]
is \emph{exact}. Suppose this point established, and let us show that it
implies the proposition. Note that $\mathfrak{K}=\Gamma(X,\mathscr{J})$ is an
ideal of square zero in $A$, and is therefore a module over
$A_0=A/\mathfrak{K}$. By hypothesis, we have $X_0=\operatorname{Spec}(A_0)$,
and since the underlying topological spaces $X_0$ and $X$ are identical,
$\mathfrak{K}=\Gamma(X_0,\mathscr{J})$; moreover, since
$\mathscr{J}^2=0$, $\mathscr{J}$ is a quasi-coherent
$(\mathscr{O}_X/\mathscr{J})$-module, so we have
$\mathscr{J}\simeq\widetilde{\mathfrak{K}}$ and
$\mathfrak{K}_x=\mathscr{J}_x$ for every $x\in X_0$
(\hyperref[I.1.4.1]{1.4.1}).

This being so, let $X'=\operatorname{Spec}(A)$, and consider the morphism of
preschemes $f=(\psi,\theta):X\to X'$ corresponding to the identity map
$A\to\Gamma(X,\mathscr{O}_X)$
(\hyperref[I.2.2.4]{2.2.4}). For every affine open subset $V$ in $X$, the
diagram
\[
  \xymatrix{
    A\ar[r]\ar[d] & \Gamma(V,\mathscr{O}_X|V)\ar[d]\\
    A_0=A/\mathfrak{K}\ar[r] & \Gamma(V,\mathscr{O}_{X_0}|V)
  }
\]
is commutative, from which we conclude that the diagram
\[
  \xymatrix{
    X' & X\ar[l]_f\\
    X'_0\ar[u]^{j'} & X_0\ar[l]^{f_0}\ar[u]_j
  }
\]
is commutative, $X'_0$ being the closed sub-prescheme of $X'$ defined by the
quasi-coherent sheaf of ideals $\widetilde{\mathfrak{K}}$, and $j$, $j'$ the
canonical injection morphisms. But since $X_0$ is affine, $f_0$ is an
isomorphism, and since the continuous maps underlying $j$ and $j'$
are the identity maps, we see first of all that
$\psi:X\to X'$ is a homeomorphism. Moreover, the relation
$\mathfrak{K}_x=\mathscr{J}_x$ shows that the restriction of
$\theta^\sharp:\psi^*(\mathscr{O}_{X'})\to\mathscr{O}_X$ is an
\emph{isomorphism} of $\psi^*(\widetilde{\mathfrak{K}})$ onto $\mathscr{J}$;
on the other hand, by passing to quotients, $\theta^\sharp$ gives an
\emph{isomorphism}
$\psi^*(\mathscr{O}_{X'}/\widetilde{\mathfrak{K}})
\to\mathscr{O}_X/\mathscr{J}$, since $f_0$ is an isomorphism; we
conclude immediately by the five lemma (M, I, 1.1) that $\theta^\sharp$ is
itself an isomorphism, hence that $f$ is an \emph{isomorphism}, and
consequently that $X$ is affine.

Everything therefore comes down to proving the exactness of
(\hyperref[I.5.1.9.1]{5.1.9.1}), which will follow from

\oldpage[I]{131}
$H^1(X,\mathscr{J})=0$. Now
$H^1(X,\mathscr{J})=H^1(X_0,\mathscr{J})$, and we have seen that
$\mathscr{J}$ is a quasi-coherent $\mathscr{O}_{X_0}$-module. Our
assertion will therefore follow from the

\begin{lemma}[5.1.9.2]
\label{I.5.1.9.2}
If $Y$ is an affine scheme and $\mathscr{F}$ a quasi-coherent
$\mathscr{O}_Y$-module, then $H^1(Y,\mathscr{F})=0$.
\end{lemma}

This lemma will be proved in Chapter III, § 1, as a consequence of the more
general theorem that $H^i(Y,\mathscr{F})=0$ for every $i>0$. To give an
independent proof of it, observe that $H^1(Y,\mathscr{F})$ is identified with
the module $\operatorname{Ext}^1_{\mathscr{O}_Y}
(Y;\mathscr{O}_Y,\mathscr{F})$ of classes of extensions of the
$\mathscr{O}_Y$-module $\mathscr{O}_Y$ by the $\mathscr{O}_Y$-module
$\mathscr{F}$ (T, 4.2.3); it all comes down, therefore, to proving that such
an extension $\mathscr{G}$ is trivial. Now, for every $y\in Y$, there is a
neighbourhood $V$ of $y$ in $Y$ such that $\mathscr{G}|V$ is isomorphic to
$\mathscr{F}|V\oplus\mathscr{O}_Y|V$
(\hyperref[0.5.4.9]{0, 5.4.9}); we conclude that $\mathscr{G}$ is a
quasi-coherent $\mathscr{O}_Y$-module. If $A$ is the ring of $Y$, we
therefore have $\mathscr{F}=\widetilde{M}$, $\mathscr{G}=\widetilde{N}$,
where $M$ and $N$ are $A$-modules, and by hypothesis $N$ is an extension of
the $A$-module $A$ by the $A$-module $M$ (\hyperref[I.1.3.11]{1.3.11}).
Since this extension is necessarily trivial, the lemma is proved, and
consequently so is (\hyperref[I.5.1.9]{5.1.9}).

\begin{corollary}[5.1.10]
\label{I.5.1.10}
Let $X$ be a prescheme such that $\mathscr{N}_X$ is nilpotent. For $X$
to be an affine scheme, it is necessary and sufficient that $X_{\mathrm{red}}$
be one.
\end{corollary}

\subsection{Existence of a sub-prescheme with a given underlying space.}
\label{subsection:I.5.2}

\begin{proposition}[5.2.1]
\label{I.5.2.1}
For every locally closed subspace $Y$ of the underlying space of a
prescheme $X$, there exists one and only one reduced sub-prescheme of $X$
having $Y$ as underlying space.
\end{proposition}

Since uniqueness follows from (\hyperref[I.5.1.2]{5.1.2}), it remains to
prove the existence of the sub-prescheme in question.

If $X$ is affine with ring $A$, and $Y$ is closed in $X$, the proposition is
immediate: $\mathfrak{j}(Y)$ is the largest ideal
$\mathfrak{a}\subset A$ such that $V(\mathfrak{a})=Y$, and it is equal to its
radical (\hyperref[I.1.1.4]{1.1.4 (i)}), so
$A/\mathfrak{j}(Y)$ is a reduced ring.

In the general case, for every affine open subset $U\subset X$ such that
$U\cap Y$ is closed in $U$, consider the closed sub-prescheme $Y_U$ of $U$
defined by the sheaf of ideals associated to the ideal $\mathfrak{j}(U\cap Y)$
of $A(U)$, which is reduced. Let us show that, if $V$ is an affine open
subset of $X$ contained in $U$, $Y_V$ is \emph{induced} by $Y_U$ on
$V\cap Y$; now, this induced prescheme is a closed sub-prescheme of $V$ which
is reduced and has $V\cap Y$ as underlying space; the uniqueness of $Y_V$
therefore implies our assertion.

\begin{proposition}[5.2.2]
\label{I.5.2.2}
Let $X$ be a reduced prescheme, $f:X\to Y$ a morphism, and $Z$ a closed
sub-prescheme of $Y$ such that $f(X)\subset Z$; then $f$ factors
as $X\xrightarrow{g}Z\xrightarrow{j}Y$, where $j$ is the injection morphism.
\end{proposition}

It follows from the hypothesis that the closed sub-prescheme $f^{-1}(Z)$ of
$X$ has the whole of $X$ as underlying space
(\hyperref[I.4.4.1]{4.4.1}); since $X$ is reduced, this closed
sub-prescheme coincides with $X$ (\hyperref[I.5.1.2]{5.1.2}), and the
proposition therefore follows from (\hyperref[I.4.4.1]{4.4.1}).

\begin{corollary}[5.2.3]
\label{I.5.2.3}
Let $X$ be a reduced sub-prescheme of a prescheme $Y$; if $Z$ is the
reduced closed sub-prescheme of $Y$ having $\overline{X}$ as underlying
space, then $X$ is a sub-prescheme induced on an open subset of $Z$.
\end{corollary}

\oldpage[I]{132}
Indeed, there is an open subset $U$ of $Y$ such that
$X=U\cap\overline{X}$; since $X$ is a reduced sub-prescheme of $Z$ by
virtue of (\hyperref[I.5.2.2]{5.2.2}), the sub-prescheme $X$ is induced by
$Z$ on the open subspace $X$ by virtue of the uniqueness
(\hyperref[I.5.2.1]{5.2.1}).

\begin{corollary}[5.2.4]
\label{I.5.2.4}
Let $f:X\to Y$ be a morphism, and $X'$ (resp. $Y'$) a closed sub-prescheme of
$X$ (resp. $Y$) defined by a quasi-coherent sheaf of ideals $\mathscr{J}$
(resp. $\mathscr{K}$) of $\mathscr{O}_X$ (resp. $\mathscr{O}_Y$). Suppose
that $X'$ is reduced and that $f(X')\subset Y'$. Then we have
$f^*(\mathscr{K})\mathscr{O}_X\subset\mathscr{J}$.
\end{corollary}

Since the restriction of $f$ to $X'$ factors as
$X'\to Y'\to Y$ by (\hyperref[I.5.2.2]{5.2.2}), it suffices to apply
(\hyperref[I.4.4.6]{4.4.6}).

\subsection{Diagonal; graph of a morphism.}
\label{subsection:I.5.3}

\begin{env}[5.3.1]
\label{I.5.3.1}
Let $X$ be an $S$-prescheme; we call the \emph{diagonal morphism} of $X$ into
$X\times_S X$, and denote by $\Delta_{X|S}$, or $\Delta_X$, or even $\Delta$ if
no confusion is possible, the $S$-morphism $(1_X,1_X)_S$, in other
words, the unique $S$-morphism $\Delta_X$ such that
\[
  p_1\circ\Delta_X=p_2\circ\Delta_X=1_X
  \tag{5.3.1.1}\label{I.5.3.1.1}
\]
denoting by $p_1$, $p_2$ the projections of $X\times_S X$
(Definition \hyperref[I.3.2.1]{3.2.1}). If $f:T\to X$, $g:T\to Y$ are two
$S$-morphisms, we verify immediately that
\[
  (f,g)_S=(f\times_S g)\circ\Delta_{T|S}
  \tag{5.3.1.2}\label{I.5.3.1.2}
\]

The reader will observe that the preceding definition and the results stated
in n\textsuperscript{os} (\hyperref[I.5.3.1]{5.3.1}) to (5.3.8) are
valid in any category, \emph{provided that the products that appear in them
exist in that category}.
\end{env}

\begin{proposition}[5.3.2]
\label{I.5.3.2}
Let $X$, $Y$ be two $S$-preschemes; if we canonically identify the product
$(X\times Y)\times(X\times Y)$ with $(X\times X)\times(Y\times Y)$, the morphism
$\Delta_{X\times Y}$ is identified with $\Delta_X\times\Delta_Y$.
\end{proposition}

Indeed, if $p_1$, $q_1$ are the first projections
$X\times X\to X$, $Y\times Y\to Y$, the first projection
$(X\times Y)\times(X\times Y)\to X\times Y$ is identified with $p_1\times q_1$,
and we have
\[
  (p_1\times q_1)\circ(\Delta_X\times\Delta_Y)
  =(p_1\circ\Delta_X)\times(q_1\circ\Delta_Y)=1_{X\times Y}
\]
the same reasoning applies to the second projections.

\begin{corollary}[5.3.4]
\label{I.5.3.4}
For every extension $S'\to S$ of the base prescheme,
$\Delta_{X_{(S')}}$ is canonically identified with $(\Delta_X)_{(S')}$.
\end{corollary}

It suffices to note that $(X\times_S X)_{(S')}$ is canonically identified with
$X_{(S')}\times_{S'}X_{(S')}$
(\hyperref[I.3.3.10]{3.3.10}).

\begin{proposition}[5.3.5]
\label{I.5.3.5}
Let $X$, $Y$ be two $S$-preschemes and $\varphi:S\to T$ a morphism, making
every $S$-prescheme a $T$-prescheme. Let $f:X\to S$, $g:Y\to S$ be the
structural morphisms, $p$, $q$ the projections of $X\times_S Y$, and
$\pi=f\circ p=g\circ q$ the structural morphism $X\times_S Y\to S$. Then the
diagram
\[
  \xymatrix{
    X\times_S Y\ar[r]^{(p,q)_T}\ar[d]_\pi &
    X\times_T Y\ar[d]^{f\times_T g}\\
    S\ar[r]_{\Delta_{S|T}} &
    S\times_T S
  }
  \tag{5.3.5.1}\label{I.5.3.5.1}
\]
\oldpage[I]{133}
is commutative, and identifies $X\times_S Y$ with the product of the
$(S\times_T S)$-preschemes $S$ and $X\times_T Y$, the projections
being identified with $\pi$ and $(p,q)_T$.
\end{proposition}

By virtue of (\hyperref[I.3.4.3]{3.4.3}), we are reduced to proving the
corresponding proposition \emph{in the category of sets}, by
replacing $X$, $Y$, $S$ by $X(Z)_T$, $Y(Z)_T$, $S(Z)_T$, $Z$ being an
arbitrary $T$-prescheme. But for the category of sets, the verification
is immediate and is left to the reader.

\begin{corollary}[5.3.6]
\label{I.5.3.6}
The morphism $(p,q)_T$ is identified (setting $P=S\times_T S$) with
$1_{X\times_T Y}\times_P\Delta_S$.
\end{corollary}

This follows from (\hyperref[I.5.3.5]{5.3.5}) and
(\hyperref[I.3.3.4]{3.3.4}).

\begin{corollary}[5.3.7]
\label{I.5.3.7}
If $f:X\to Y$ is an $S$-morphism, the diagram
\[
  \xymatrix{
    X\ar[r]^{(1_X,f)_S}\ar[d]_f &
    X\times_S Y\ar[d]^{f\times_S 1_Y}\\
    Y\ar[r]_{\Delta_Y} &
    Y\times_S Y
  }
\]
is commutative, and identifies $X$ with the product of the
$(Y\times_S Y)$-preschemes $Y$ and $X\times_S Y$.
\end{corollary}

It suffices to apply (\hyperref[I.5.3.5]{5.3.5}), replacing $S$ by
$Y$ and $T$ by $S$, and noting that $X\times_Y Y=X$
(\hyperref[I.3.3.3]{3.3.3}).

\begin{proposition}[5.3.8]
\label{I.5.3.8}
For $f:X\to Y$ to be a monomorphism of preschemes, it is necessary and
sufficient that $\Delta_{X|Y}$ be an isomorphism of $X$ onto $X\times_Y X$.
\end{proposition}

Indeed, to say that $f$ is a monomorphism means that for every
$Y$-prescheme $Z$, the corresponding map
$f':X(Z)_Y\to Y(Z)_Y$ is an injection, and since $Y(Z)_Y$ consists of a single
element, this means that $X(Z)_Y$ has at most one element. But this can also
be expressed by saying that $X(Z)_Y\times X(Z)_Y$ is canonically isomorphic to
$X(Z)_Y$, and since the first of these sets is $(X\times_Y X)(Z)_Y$
(\hyperref[I.3.4.3.1]{3.4.3.1}), this means that $\Delta_{X|Y}$ is an
isomorphism.

\begin{proposition}[5.3.9]
\label{I.5.3.9}
The diagonal morphism $\Delta_X$ is an immersion of $X$ into
$X\times_S X$.
\end{proposition}

To obtain a correct proof, it
suffices to use (\hyperref[I.4.2.4]{4.2.4, $a$}): for every $x\in X$, choose an affine open neighbourhood $V\subset S$
of its image and an affine open neighbourhood $U$ of $x$ in $X$ whose
image is contained in $V$. Then $U\times_SU=U\times_VU$ is an
affine open neighbourhood of $\Delta_X(x)$;
taking into account (\hyperref[I.5.3.16]{5.3.16}) (whose proof uses only the definition
(\hyperref[I.5.3.1.1]{5.3.1.1})), we are reduced to proving (\hyperref[I.5.3.9]{5.3.9}) when
$S=\operatorname{Spec}(B)$ and $X=\operatorname{Spec}(A)$ are affine
schemes; it is then clear by virtue of (\hyperref[I.5.3.1.1]{5.3.1.1}) that $\Delta_X$ corresponds to
the canonical homomorphism $A\otimes_BA\to A$ which sends $x\otimes y$ to
$xy$; this homomorphism being surjective, $\Delta_X$ is in this case a
closed immersion (\hyperref[I.4.2.3]{4.2.3}), which completes the proof.

We say that the sub-prescheme of $X\times_S X$ associated to the immersion
$\Delta_X$ (\hyperref[I.4.2.1]{4.2.1}) is \emph{the diagonal} of
$X\times_S X$.

\begin{corollary}[5.3.10]
\label{I.5.3.10}
Under the hypotheses of (\hyperref[I.5.3.5]{5.3.5}), $(p,q)_T$ is an
immersion.
\end{corollary}

This follows from (\hyperref[I.5.3.6]{5.3.6}) and
(\hyperref[I.4.3.1]{4.3.1}).

We say (under the hypotheses of (\hyperref[I.5.3.5]{5.3.5})) that $(p,q)_T$
is the \emph{canonical immersion} of $X\times_S Y$ into $X\times_T Y$.

\begin{corollary}[5.3.11]
\label{I.5.3.11}
Let $X$, $Y$ be two $S$-preschemes and $f:X\to Y$ an $S$-morphism; then the
graph morphism $\Gamma_f=(1_X,f)_S$ of $f$
(\hyperref[I.3.3.14]{3.3.14}) is an immersion of $X$ into
$X\times_S Y$.
\end{corollary}

This is the particular case of Corollary (\hyperref[I.5.3.10]{5.3.10})
where we replace $S$ by $Y$ and $T$ by $S$
(cf. (\hyperref[I.5.3.7]{5.3.7})).

\oldpage[I]{134}
The sub-prescheme of $X\times_S Y$ associated to the immersion $\Gamma_f$
(\hyperref[I.4.2.1]{4.2.1}) is called the \emph{graph} of the morphism $f$;
the sub-preschemes of $X\times_S Y$ that are graphs of morphisms
$X\to Y$ are characterized by the fact that the restriction to such a
sub-prescheme $G$ of the projection $p_1:X\times_S Y\to X$ is an
\emph{isomorphism} $g$ of $G$ onto $X$: $G$ is then the graph of the morphism
$p_2\circ g^{-1}$, where $p_2$ is the projection $X\times_S Y\to Y$.

When in particular we take $X=S$, the $S$-morphisms $S\to Y$, which are
none other than the $S$-sections of $Y$ (\hyperref[I.2.5.5]{2.5.5}), are
equal to their graph morphisms; the sub-preschemes of $Y$ that are
graphs of $S$-sections (in other words, those that are isomorphic to $S$ via
the restriction of the structural morphism $Y\to S$) are also called the
\emph{images} of these sections, or, by abuse of language, the
\emph{$S$-sections} of $Y$.

\begin{corollary}[5.3.12]
\label{I.5.3.12}
The hypotheses and notation being those of
(\hyperref[I.5.3.11]{5.3.11}), for every morphism $g:S'\to S$, let $f'$
be the inverse image of $f$ by $g$ (\hyperref[I.3.3.7]{3.3.7}); then
$\Gamma_{f'}$ is the inverse image of $\Gamma_f$ by $g$.
\end{corollary}

This is a particular case of formula
(\hyperref[I.3.3.10.1]{3.3.10.1}).

\begin{corollary}[5.3.13]
\label{I.5.3.13}
Let $f:X\to Y$, $g:Y\to Z$ be two morphisms; if $g\circ f$ is an
immersion (resp. a local immersion), then so is $f$.
\end{corollary}

Indeed, $f$ factors as
\[
  X\xrightarrow{\Gamma_f}X\times_Z Y\xrightarrow{p_2}Y.
\]
On the other hand, $p_2$ is identified with $(g\circ f)\times_Z 1_Y$
(\hyperref[I.3.3.4]{3.3.4}); if $g\circ f$ is an immersion (resp. a
local immersion), then so is $p_2$
(\hyperref[I.4.3.1]{4.3.1} and \hyperref[I.4.5.5]{4.5.5}), and since
$\Gamma_f$ is an immersion (\hyperref[I.5.3.11]{5.3.11}), we conclude
by (\hyperref[I.4.2.5]{4.2.5}) (resp.
(\hyperref[I.4.5.5]{4.5.5})).

\begin{corollary}[5.3.14]
\label{I.5.3.14}
Let $j:X\to Y$, $g:X\to Z$ be two $S$-morphisms. If $j$ is an immersion
(resp. a local immersion), then so is $(j,g)_S$.
\end{corollary}

Indeed, if $p:Y\times_S Z\to Y$ is the first projection, we have
$j=p\circ(j,g)_S$, and it suffices to apply
(\hyperref[I.5.3.13]{5.3.13}).

\begin{proposition}[5.3.15]
\label{I.5.3.15}
If $f:X\to Y$ is an $S$-morphism, the diagram
\[
  \xymatrix{
    X\ar[r]^{\Delta_X}\ar[d]_f &
    X\times_S X\ar[d]^{f\times_S f}\\
    Y\ar[r]_{\Delta_Y} &
    Y\times_S Y
  }
  \tag{5.3.15.1}\label{I.5.3.15.1}
\]
is commutative (in other words, $\Delta_X$ is a functorial morphism
in the category of preschemes).
\end{proposition}

The verification is immediate and is left to the reader.

\begin{corollary}[5.3.16]
\label{I.5.3.16}
If $X$ is a sub-prescheme of $Y$, the diagonal $\Delta_X(X)$ is identified with
a sub-prescheme of $\Delta_Y(Y)$, whose underlying space is identified with
\[
  \Delta_Y(Y)\cap p_1^{-1}(X)=\Delta_Y(Y)\cap p_2^{-1}(X)
\]
($p_1$, $p_2$ being the projections of $Y\times_S Y$).
\end{corollary}

Let us apply (\hyperref[I.5.3.15]{5.3.15}) to the injection morphism
$f:X\to Y$; we then know that $f\times_S f$ is an immersion, identifying
the underlying space of $X\times_S X$ with the subspace
$p_1^{-1}(X)\cap p_2^{-1}(X)$ of $Y\times_S Y$
(\hyperref[I.4.3.1]{4.3.1}); moreover, if
$z\in\Delta_Y(Y)\cap p_1^{-1}(X)$, we have $z=\Delta_Y(y)$

\oldpage[I]{135}
and $y=p_1(z)\in X$, so $y=f(y)$, and $z=\Delta_Y(f(y))$ belongs to
$\Delta_X(X)$ by virtue of the commutativity of the diagram
(\hyperref[I.5.3.15.1]{5.3.15.1}).

\begin{corollary}[5.3.17]
\label{I.5.3.17}
Let $f_1:Y\to X$, $f_2:Y\to X$ be two $S$-morphisms, and $y$ a point of $Y$
such that $f_1(y)=f_2(y)=x$ and such that the homomorphisms $k(x)\to k(y)$
corresponding to $f_1$ and $f_2$ are identical. Then, if
$f=(f_1,f_2)_S$, the point $f(y)$ belongs to the diagonal
$\Delta_{X|S}(X)$.
\end{corollary}

The two homomorphisms $k(x)\to k(y)$ corresponding to $f_i$ ($i=1,2$)
define two $S$-morphisms
$g_i:\operatorname{Spec}(k(y))\to\operatorname{Spec}(k(x))$ such that the
diagrams
\[
  \xymatrix{
    \operatorname{Spec}(k(y))\ar[r]^{g_i}\ar[d] &
    \operatorname{Spec}(k(x))\ar[d]\\
    Y\ar[r]_{f_i} & X
  }
\]
are commutative. The diagram
\[
  \xymatrix{
    \operatorname{Spec}(k(y))\ar[rr]^{(g_1,g_2)_S}\ar[d] & &
    \operatorname{Spec}(k(x))\times_S\operatorname{Spec}(k(x))\ar[d]\\
    Y\ar[rr]_{(f_1,f_2)_S} & & X\times_S X
  }
\]
is therefore also commutative. Now it follows from the equality $g_1=g_2$
that the image under $(g_1,g_2)_S$ of the unique point of
$\operatorname{Spec}(k(y))$ belongs to the diagonal of
$\operatorname{Spec}(k(x))\times_S\operatorname{Spec}(k(x))$; the conclusion
therefore follows from (\hyperref[I.5.3.15]{5.3.15}).

\subsection{Separated morphisms and preschemes.}
\label{subsection:I.5.4}

\begin{definition}[5.4.1]
\label{I.5.4.1}
We say that a morphism of preschemes $f:X\to Y$ is \emph{separated} if the
diagonal morphism $X\to X\times_Y X$ is a \emph{closed immersion}; we
then also say that $X$ is a \emph{prescheme separated over $Y$}, or a
\emph{$Y$-scheme}. We say that a prescheme $X$ is separated if it is
separated over $\operatorname{Spec}(\mathbf{Z})$; we then also say that $X$
is a \emph{scheme} (cf. (\hyperref[I.5.5.7]{5.5.7})).
\end{definition}

By virtue of (\hyperref[I.5.3.9]{5.3.9}), for $X$ to be separated over
$Y$, it is necessary and sufficient that $\Delta_X(X)$ be a \emph{closed
subspace} of the underlying space of $X\times_Y X$.

\begin{proposition}[5.4.2]
\label{I.5.4.2}
Let $S\to T$ be a separated morphism. If $X$ and $Y$ are two $S$-preschemes,
the canonical immersion $X\times_S Y\to X\times_T Y$
(\hyperref[I.5.3.10]{5.3.10}) is closed.
\end{proposition}

Indeed, referring to the diagram
(\hyperref[I.5.3.5.1]{5.3.5.1}), we see that $(p,q)_T$ can be regarded
as obtained from $\Delta_{S|T}$ by the extension
$f\times_T g:X\times_T Y\to S\times_T S$ of the base prescheme $S\times_T S$;
the proposition then follows from (\hyperref[I.4.3.2]{4.3.2}).

\begin{corollary}[5.4.3]
\label{I.5.4.3}
Let $Y$ be an $S$-scheme and $f:X\to Y$ an $S$-morphism. Then the graph
morphism $\Gamma_f:X\to X\times_S Y$ (\hyperref[I.5.3.11]{5.3.11}) is a
closed immersion.
\end{corollary}

This is the particular case of (\hyperref[I.5.4.2]{5.4.2}) where we replace
$S$ by $Y$ and $T$ by $S$.

\begin{corollary}[5.4.4]
\label{I.5.4.4}
Let $f:X\to Y$, $g:Y\to Z$ be two morphisms, $g$ being separated. If
$g\circ f$ is a closed immersion, then so is $f$.
\end{corollary}

The proof from (\hyperref[I.5.4.3]{5.4.3}) is the same as
that of (\hyperref[I.5.3.13]{5.3.13}) from
(\hyperref[I.5.3.11]{5.3.11}).

\oldpage[I]{136}
\begin{corollary}[5.4.5]
\label{I.5.4.5}
Let $Z$ be an $S$-scheme and $j:X\to Y$, $g:X\to Z$ two $S$-morphisms. If $j$
is a closed immersion, then so is
$(j,g)_S:X\to Y\times_S Z$.
\end{corollary}

The proof from (\hyperref[I.5.4.4]{5.4.4}) is the same as
that of (\hyperref[I.5.3.14]{5.3.14}) from
(\hyperref[I.5.3.13]{5.3.13}).

\begin{corollary}[5.4.6]
\label{I.5.4.6}
If $X$ is an $S$-scheme, every $S$-section of $X$
(\hyperref[I.2.5.5]{2.5.5}) is a closed immersion.
\end{corollary}

If $\varphi:X\to S$ is the structural morphism and $\psi:S\to X$ an
$S$-section of $X$, it suffices to apply
(\hyperref[I.5.4.5]{5.4.5}) to $\varphi\circ\psi=1_S$.

\begin{corollary}[5.4.7]
\label{I.5.4.7}
Let $S$ be an integral prescheme, $s$ its generic point, and $X$ an
$S$-scheme. If two $S$-sections $f$, $g$ of $X$ are such that $f(s)=g(s)$,
then $f=g$.
\end{corollary}

Indeed, if $x=f(s)=g(s)$, the homomorphisms $k(x)\to k(s)$ corresponding to
$f$ and $g$ are necessarily identical. If $h=(f,g)_S$, we deduce from this
(\hyperref[I.5.3.17]{5.3.17}) that $h(s)$ belongs to the diagonal
$Z=\Delta_X(X)$; but since $S=\overline{\{s\}}$ and $Z$ is closed by
hypothesis, we have $h(S)\subset Z$. It then follows from
(\hyperref[I.5.2.2]{5.2.2}) that $h$ factors as
$S\to Z\to X\times_S X$, and we conclude that $f=g$ by the definition of the
diagonal.

\begin{remark}[5.4.8]
\label{I.5.4.8}
If we suppose conversely that the conclusion of
(\hyperref[I.5.4.3]{5.4.3}) holds when $f=1_Y$, we conclude
that $Y$ is separated over $S$; similarly, if we suppose that the conclusion
of (\hyperref[I.5.4.4]{5.4.4}) applies to the two morphisms
\[
  Y\xrightarrow{\Delta_Y}Y\times_Z Y\xrightarrow{p_1}Y,
\]
we deduce that $\Delta_Y$ is a closed immersion, hence that $Y$ is separated
over $Z$; finally, the validity of the conclusion of
(\hyperref[I.5.4.6]{5.4.6}) for the $Y$-section $\Delta_Y$ of the
$Y$-prescheme $Y\times_S Y\to Y$ implies that $Y$ is separated over
$S$.
\end{remark}


\subsection{Criteria for separatedness.}
\label{subsection:I.5.5}

\begin{proposition}[5.5.1]
\label{I.5.5.1}
\begin{enumerate}
  \item[\rm(i)] Every monomorphism of preschemes (and in particular every
    immersion) is a separated morphism.
  \item[\rm(ii)] The composite of two separated morphisms is separated.
  \item[\rm(iii)] If $f:X\to X'$, $g:Y\to Y'$ are two separated
    $S$-morphisms, $f\times_S g$ is separated.
  \item[\rm(iv)] If $f:X\to Y$ is a separated $S$-morphism, the
    $S'$-morphism $f_{(S')}$ is separated for every extension $S'\to S$ of the
    base prescheme.
  \item[\rm(v)] If the composite $g\circ f$ of two morphisms is separated, $f$
    is separated.
  \item[\rm(vi)] For a morphism $f$ to be separated, it is necessary and
    sufficient that $f_{\mathrm{red}}$ (\hyperref[I.5.1.5]{5.1.5}) be
    separated.
\end{enumerate}
\end{proposition}

(i) follows immediately from (\hyperref[I.5.3.8]{5.3.8}). If $f:X\to Y$,
$g:Y\to Z$ are two morphisms, the diagram
\[
  \xymatrix{
    X\ar[rr]^{\Delta_{X|Z}}\ar[dr]_{\Delta_{X|Y}} & &
    X\times_Z X\\
    & X\times_Y X\ar[ur]_j
  }
  \tag{5.5.1.1}\label{I.5.5.1.1}
\]
where $j$ denotes the canonical immersion (\hyperref[I.5.3.10]{5.3.10}), is
commutative, as we verify immediately. If $f$ and $g$ are separated,
$\Delta_{X|Y}$ is a closed immersion by definition, and $j$ is a
closed immersion by virtue of (\hyperref[I.5.4.2]{5.4.2}), so
$\Delta_{X|Z}$ is a closed immersion by
(\hyperref[I.4.2.4]{4.2.4}), which

\oldpage[I]{137}
proves (ii). In view of (i) and (ii), (iii) and (iv) are equivalent
(\hyperref[I.3.5.1]{3.5.1}), and it suffices to prove (iv). Now
$X_{(S')}\times_{Y_{(S')}}X_{(S')}$ is canonically identified with
$(X\times_Y X)\times_Y Y_{(S')}$ by virtue of
(\hyperref[I.3.3.11]{3.3.11}) and
(\hyperref[I.3.3.9.1]{3.3.9.1}), and we verify immediately that the
diagonal morphism $\Delta_{X_{(S')}}$ is then identified with
$\Delta_X\times_Y 1_{Y_{(S')}}$; the proposition therefore follows from
(\hyperref[I.4.3.1]{4.3.1}).

To establish (v), consider, as in
(\hyperref[I.5.3.13]{5.3.13}), the factorization
$X\xrightarrow{\Gamma_f}X\times_Z Y\xrightarrow{p_2}Y$ of $f$, noting
that $p_2=(g\circ f)\times_Z 1_Y$; the hypothesis that $g\circ f$ is separated
implies that $p_2$ is separated by (iii) and (i), and since $\Gamma_f$ is an
immersion, $\Gamma_f$ is separated by (i), so $f$ is separated by (ii).
Finally, to prove (vi), recall that the preschemes
$X_{\mathrm{red}}\times_{Y_{\mathrm{red}}}X_{\mathrm{red}}$ and
$X_{\mathrm{red}}\times_Y X_{\mathrm{red}}$ are canonically identified
(\hyperref[I.5.1.7]{5.1.7}); if we denote by $j$ the injection
$X_{\mathrm{red}}\to X$, the diagram
\[
  \xymatrix{
    X_{\mathrm{red}}\ar[r]^{\Delta_{X_{\mathrm{red}}}}\ar[d]_j &
    X_{\mathrm{red}}\times_Y X_{\mathrm{red}}\ar[d]^{j\times_Y j}\\
    X\ar[r]_{\Delta_X} & X\times_Y X
  }
\]
is commutative (\hyperref[I.5.3.15]{5.3.15}), and the proposition follows
from the fact that the vertical arrows are homeomorphisms of the underlying
spaces (\hyperref[I.4.3.1]{4.3.1}).

\begin{corollary}[5.5.2]
\label{I.5.5.2}
If $f:X\to Y$ is separated, the restriction of $f$ to every sub-prescheme of
$X$ is separated.
\end{corollary}

This follows from (\hyperref[I.5.5.1]{5.5.1, (i) and (ii)}).

\begin{corollary}[5.5.3]
\label{I.5.5.3}
If $X$, $Y$ are two $S$-preschemes such that $Y$ is separated over
$S$, then $X\times_S Y$ is separated over $X$.
\end{corollary}

This is a particular case of (\hyperref[I.5.5.1]{5.5.1, (iv)}).

\begin{proposition}[5.5.4]
\label{I.5.5.4}
Let $X$ be a prescheme, and suppose that its underlying space is the union of
a finite family of closed subsets $X_k$ $(1\leq k\leq n)$; for each $k$ we
consider the reduced sub-prescheme of $X$ having $X_k$ as underlying space
(\hyperref[I.5.2.1]{5.2.1}), and we again denote it by $X_k$. Let
$f:X\to Y$ be a morphism, and for each $k$, let $Y_k$ be a closed subset of
$Y$ such that $f(X_k)\subset Y_k$; we again denote by $Y_k$ the reduced
sub-prescheme of $Y$ having $Y_k$ as underlying space, so that the restriction
$X_k\to Y$ of $f$ to $X_k$ factors as
$X_k\xrightarrow{f_k}Y_k\to Y$ (\hyperref[I.5.2.2]{5.2.2}). For $f$
to be separated, it is necessary and sufficient that the $f_k$ be separated.
\end{proposition}

Necessity follows from
(\hyperref[I.5.5.1]{5.5.1, (i), (ii) and (v)}). Conversely, if the
condition of the statement is satisfied, each of the restrictions $X_k\to Y$
of $f$ is separated (\hyperref[I.5.5.1]{5.5.1, (i) and (ii)}); if $p_1$,
$p_2$ are the projections of $X\times_Y X$, the subspace
$\Delta_{X_k}(X_k)$ is identified with the subspace
$\Delta_X(X)\cap p_1^{-1}(X_k)$ of the underlying space of
$X\times_Y X$ (\hyperref[I.5.3.16]{5.3.16}); since these subspaces are
closed in $X\times_Y X$, so is their union $\Delta_X(X)$.

Suppose in particular that the $X_k$ are the
\emph{irreducible components} of $X$; then we may suppose that the $Y_k$
are irreducible components of $Y$
(\hyperref[0.2.1.5]{0, 2.1.5}); in this case Proposition
(\hyperref[I.5.5.4]{5.5.4}) therefore reduces the notion of separatedness
to the case of \emph{integral} preschemes (\hyperref[I.2.1.7]{2.1.7}).

\oldpage[I]{138}

\begin{proposition}[5.5.5]
\label{I.5.5.5}
Let $(Y_\lambda)$ be an open covering of a prescheme $Y$; for a
morphism $f:X\to Y$ to be separated, it is necessary and sufficient that each
of its restrictions $f^{-1}(Y_\lambda)\to Y_\lambda$ be separated.
\end{proposition}

If we set $X_\lambda=f^{-1}(Y_\lambda)$, it all comes down, taking into
account (\hyperref[I.4.2.4]{4.2.4, b)}) and the identity of the products
$X_\lambda\times_Y X_\lambda$ and
$X_\lambda\times_{Y_\lambda}X_\lambda$
(\hyperref[I.3.2.5]{3.2.5}), to proving that the
$X_\lambda\times_Y X_\lambda$ form a covering of $X\times_Y X$. Now,
if we set $Y_{\lambda\mu}=Y_\lambda\cap Y_\mu$ and
$X_{\lambda\mu}=X_\lambda\cap X_\mu=f^{-1}(Y_{\lambda\mu})$,
$X_\lambda\times_Y X_\mu$ is identified with the product
$X_{\lambda\mu}\times_{Y_{\lambda\mu}}X_{\lambda\mu}$
(\hyperref[I.3.2.6.4]{3.2.6.4}), hence also with
$X_{\lambda\mu}\times_Y X_{\lambda\mu}$
(\hyperref[I.3.2.5]{3.2.5}), and finally with an open subset of
$X_\lambda\times_Y X_\lambda$, which establishes our assertion
(\hyperref[I.3.2.7]{3.2.7}).

Proposition (\hyperref[I.5.5.4]{5.5.4}) makes it possible, by taking a
covering of $Y$ by affine open subsets, to reduce the study of separated
morphisms to that of separated morphisms with values in affine schemes.

\begin{proposition}[5.5.6]
\label{I.5.5.6}
Let $Y$ be an affine scheme, $X$ a prescheme, and $(U_\alpha)$ a covering of
$X$ by affine open subsets. For a morphism $f:X\to Y$ to be separated, it
is necessary and sufficient that, for every pair of indices $(\alpha,\beta)$,
$U_\alpha\cap U_\beta$ be an affine open subset, and that the ring
$\Gamma(U_\alpha\cap U_\beta,\mathscr{O}_X)$ be generated by the union of the
canonical images of the rings $\Gamma(U_\alpha,\mathscr{O}_X)$ and
$\Gamma(U_\beta,\mathscr{O}_X)$.
\end{proposition}

The $U_\alpha\times_Y U_\beta$ form an open covering of
$X\times_Y X$ (\hyperref[I.3.2.7]{3.2.7}); denoting by $p$ and $q$ the
projections of $X\times_Y X$, we have
\[
  \begin{aligned}
    \Delta_X^{-1}(U_\alpha\times_Y U_\beta)
      &=\Delta_X^{-1}\bigl(p^{-1}(U_\alpha)\cap q^{-1}(U_\beta)\bigr)\\
      &=\Delta_X^{-1}\bigl(p^{-1}(U_\alpha)\bigr)
        \cap\Delta_X^{-1}\bigl(q^{-1}(U_\beta)\bigr)
        =U_\alpha\cap U_\beta~;
  \end{aligned}
\]
it all comes down, therefore, to expressing that the restriction of
$\Delta_X$ to $U_\alpha\cap U_\beta$ is a closed immersion into
$U_\alpha\times_Y U_\beta$. Now this restriction is none other than
$(j_\alpha,j_\beta)_Y$, denoting by $j_\alpha$ (resp. $j_\beta$) the
injection morphism of $U_\alpha\cap U_\beta$ into $U_\alpha$ (resp.
$U_\beta$), as follows from the definitions. Since
$U_\alpha\times_Y U_\beta$ is an affine scheme whose ring is
canonically isomorphic to
$\Gamma(U_\alpha,\mathscr{O}_X)
 \otimes_{\Gamma(Y,\mathscr{O}_Y)}
 \Gamma(U_\beta,\mathscr{O}_X)$
(\hyperref[I.3.2.2]{3.2.2}), we see that $U_\alpha\cap U_\beta$ must be
an affine scheme and that the map
$h_\alpha\otimes h_\beta\to h_\alpha h_\beta$ from the ring
$A(U_\alpha\times_Y U_\beta)$ to
$\Gamma(U_\alpha\cap U_\beta,\mathscr{O}_X)$ must be surjective
(\hyperref[I.4.2.3]{4.2.3}), which completes the proof.

\begin{corollary}[5.5.7]
\label{I.5.5.7}
An affine scheme is separated (and is consequently a scheme, which justifies
the terminology of (\hyperref[I.5.4.1]{5.4.1})).
\end{corollary}

\begin{corollary}[5.5.8]
\label{I.5.5.8}
Let $Y$ be an affine scheme; for $f:X\to Y$ to be a separated morphism, it
is necessary and sufficient that $X$ be separated (in other words, that $X$
be a scheme).
\end{corollary}

Indeed, we observe that the criterion of
(\hyperref[I.5.5.6]{5.5.6}) does not depend on $f$.

\begin{corollary}[5.5.9]
\label{I.5.5.9}
For a morphism $f:X\to Y$ to be separated, it is necessary that for every
open subset $U$ on which $Y$ induces a separated prescheme, the induced
prescheme $f^{-1}(U)$ be separated, and it is sufficient that this be so for
every affine open subset $U\subset Y$.
\end{corollary}

The necessity of the condition follows from
(\hyperref[I.5.5.4]{5.5.4}) and
(\hyperref[I.5.5.1]{5.5.1, (ii)}); the sufficiency follows from
(\hyperref[I.5.5.4]{5.5.4}) and
(\hyperref[I.5.5.8]{5.5.8}), taking into account the existence of affine
open coverings of $Y$.

In particular, if $X$ and $Y$ are affine schemes, \emph{every} morphism
$X\to Y$ is separated.

\begin{proposition}[5.5.10]
\label{I.5.5.10}
Let $Y$ be a scheme and $f:X\to Y$ a morphism. For every affine open subset
$U$ of $X$ and every affine open subset $V$ of $Y$, $U\cap f^{-1}(V)$ is
affine.
\end{proposition}

Let $p_1$, $p_2$ be the projections of $X\times_{\mathbf{Z}}Y$; the
subspace $U\cap f^{-1}(V)$ is the image under $p_1$ of
$\Gamma_f(X)\cap p_1^{-1}(U)\cap p_2^{-1}(V)$. Now
$p_1^{-1}(U)\cap p_2^{-1}(V)$ is identified with the underlying space of the

\oldpage[I]{139}
prescheme $U\times_{\mathbf{Z}}V$
(\hyperref[I.3.2.7]{3.2.7}), and is consequently an affine scheme
(\hyperref[I.3.2.2]{3.2.2}); since $\Gamma_f(X)$ is closed in
$X\times_{\mathbf{Z}}Y$ (\hyperref[I.5.4.3]{5.4.3}),
$\Gamma_f(X)\cap p_1^{-1}(U)\cap p_2^{-1}(V)$ is closed in
$U\times_{\mathbf{Z}}V$, and consequently the prescheme induced by the
sub-prescheme of $X\times_{\mathbf{Z}}Y$ associated to $\Gamma_f$
(\hyperref[I.4.2.1]{4.2.1}) on the open subset
$\Gamma_f(X)\cap p_1^{-1}(U)\cap p_2^{-1}(V)$ of its underlying space is
a closed sub-prescheme of an affine scheme, hence an affine scheme
(\hyperref[I.4.2.3]{4.2.3}). The proposition then follows from the fact
that $\Gamma_f$ is an immersion.

\begin{examples}[5.5.11]
\label{I.5.5.11}
The prescheme of Example (\hyperref[I.2.3.2]{2.3.2})
(``projective line over a field $K$'') is \emph{separated}, because for the
covering $(X_1,X_2)$ of $X$ by affine open subsets,
$X_1\cap X_2=U_{12}$ is affine and $\Gamma(U_{12},\mathscr{O}_X)$, the ring of
rational fractions of the form $f(s)/s^m$ with $f\in K[s]$, is generated
by $K[s]$ and by $1/s$, so the conditions of
(\hyperref[I.5.5.6]{5.5.6}) are satisfied.

With the same choice of $X_1$, $X_2$, $U_{12}$ and $U_{21}$ as in Example
(\hyperref[I.2.3.2]{2.3.2}), let us this time take for $u_{12}$
the isomorphism which sends $f(s)$ to $f(t)$; this time we obtain
by gluing an \emph{integral non-separated} prescheme $X$, because the first
condition of (\hyperref[I.5.5.6]{5.5.6}) is satisfied, but not the second.
It is immediate here that
$\Gamma(X,\mathscr{O}_X)\to\Gamma(X_1,\mathscr{O}_X)=K[s]$ is an
isomorphism; the inverse isomorphism defines a morphism
$f:X\to\operatorname{Spec}(K[s])$ which is surjective, and for every
$y\in\operatorname{Spec}(K[s])$ such that $\mathfrak{j}_y\neq(0)$,
$f^{-1}(y)$ consists of a single point, but for $\mathfrak{j}_y=(0)$,
$f^{-1}(y)$ consists of \emph{two} distinct points (we say that $X$ is the
``affine line over $K$, with the point $0$ doubled'').

We can also give examples where \emph{neither} of the two conditions of
(\hyperref[I.5.5.6]{5.5.6}) is satisfied. Let us first note that in the
prime spectrum $Y$ of the polynomial ring $A=K[s,t]$ in two indeterminates
over a field $K$, the open subset $U$, the union of $D(s)$ and $D(t)$,
\emph{is not an affine open subset}. Indeed, if $z$ is a section of
$\mathscr{O}_Y$ over $U$, there exist two integers $m\geq0$, $n\geq0$ such
that $s^m z$ and $t^n z$ are the restrictions to $U$ of polynomials in $s$
and $t$ (\hyperref[I.1.4.1]{1.4.1}), which is evidently possible only if
the section $z$ extends to a section over the whole of $Y$,
identified with a polynomial in $s$ and $t$. If $U$ were an affine open
subset, the injection morphism $U\to Y$ would therefore be an isomorphism
(\hyperref[I.1.7.3]{1.7.3}), which is absurd since $U\neq Y$.

This being so, take two affine schemes $Y_1$, $Y_2$, the prime spectra of
the rings $A_1=K[s_1,t_1]$, $A_2=K[s_2,t_2]$; take
$U_{12}=D(s_1)\cup D(t_1)$, $U_{21}=D(s_2)\cup D(t_2)$, and for $u_{12}$ the
restriction to $U_{21}$ of the isomorphism $Y_2\to Y_1$ corresponding to
the ring isomorphism which sends $f(s_1,t_1)$ to $f(s_2,t_2)$;
we thus have an example where none of the conditions of
(\hyperref[I.5.5.6]{5.5.6}) is satisfied (the integral prescheme thus
obtained is called the ``affine plane over $K$, with the point $0$
doubled'').
\end{examples}

\begin{remark}[5.5.12]
\label{I.5.5.12}
\normalfont Given a property \textbf{P} of morphisms of preschemes,
consider the following propositions:
\begin{enumerate}
  \item[\rm(i)] \emph{Every closed immersion has the property \textbf{P}.}
  \item[\rm(ii)] \emph{The composite of two morphisms having the property
    \textbf{P} has the property \textbf{P}.}
  \item[\rm(iii)] \emph{If $f:X\to X'$, $g:Y\to Y'$ are two $S$-morphisms
    having the property \textbf{P}, then $f\times_S g$ has the property
    \textbf{P}.}

\oldpage[I]{140}
  \item[\rm(iv)] \emph{If $f:X\to Y$ is an $S$-morphism having the
    property \textbf{P}, every $S'$-morphism $f_{(S')}$ obtained by an
    extension $S'\to S$ of the base prescheme has the property
    \textbf{P}.}
  \item[\rm(v)] \emph{If the composite $g\circ f$ of two morphisms
    $f:X\to Y$, $g:Y\to Z$ has the property \textbf{P}, and if $g$ is
    separated, then $f$ has the property \textbf{P}.}
  \item[\rm(vi)] \emph{If a morphism $f:X\to Y$ has the
    property \textbf{P}, then so does $f_{\mathrm{red}}$
    (\hyperref[I.5.1.5]{5.1.5}).}
\end{enumerate}

Under these conditions, if we suppose (i) and (ii) satisfied, (iii) and (iv)
are \emph{equivalent}, and (v) and (vi) are \emph{consequences} of (i), (ii)
and (iii).

The first assertion has already been proved
(\hyperref[I.3.5.1]{3.5.1}). Consider the factorization
(\hyperref[I.5.3.13]{5.3.13}) of $f$ as
$X\xrightarrow{\Gamma_f}X\times_Z Y\xrightarrow{p_2}Y$; the relation
$p_2=(g\circ f)\times_Z 1_Y$ shows that if $g\circ f$ has the property
\textbf{P}, then so does $p_2$ by virtue of (iii); if $g$ is separated,
$\Gamma_f$ is a closed immersion
(\hyperref[I.5.4.3]{5.4.3}), and so also has the property \textbf{P} by
(i); finally, by virtue of (ii), $f$ has the property \textbf{P}.

Finally, consider the commutative diagram
\[
  \xymatrix{
    X_{\mathrm{red}}\ar[r]^{f_{\mathrm{red}}}\ar[d] &
    Y_{\mathrm{red}}\ar[d]\\
    X\ar[r]_f &
    Y
  }
\]
where the vertical arrows are closed immersions
(\hyperref[I.5.1.5]{5.1.5}), and therefore have the property \textbf{P} by
(i). The hypothesis that $f$ has the property \textbf{P} therefore implies by
(ii) that
$X_{\mathrm{red}}\xrightarrow{f_{\mathrm{red}}}Y_{\mathrm{red}}\to Y$
has the property \textbf{P}; finally, since a closed immersion is
separated (\hyperref[I.5.5.1]{5.5.1, (i)}), $f_{\mathrm{red}}$ has the
property \textbf{P} by virtue of (v).

Note that if we consider the propositions:
\begin{enumerate}
  \item[\rm(i')] \emph{Every immersion has the property \textbf{P}.}
  \item[\rm(v')] \emph{If $g\circ f$ has the property \textbf{P}, then so
    does $f$;}
\end{enumerate}
then the reasoning above shows that (v') is a consequence of (i'),
(ii) and (iii).
\end{remark}

\begin{env}[5.5.13]
\label{I.5.5.13}
Note that (v) and (vi) are also consequences of (i), (iii) and
\begin{enumerate}
  \item[\rm(ii')] \emph{If $j:X\to Y$ is a closed immersion and
    $g:Y\to Z$ a morphism having the property \textbf{P}, then
    $g\circ j$ has the property \textbf{P}.}
\end{enumerate}
Similarly, (v') is a consequence of (i'), (iii) and
\begin{enumerate}
  \item[\rm(ii'')] \emph{If $j:X\to Y$ is an immersion and $g:Y\to Z$ a
    morphism having the property \textbf{P}, then $g\circ j$ has the
    property \textbf{P}.}
\end{enumerate}
Indeed, this follows immediately from the arguments of
(\hyperref[I.5.5.12]{5.5.12}).
\end{env}
